\documentclass[11pt]{report}
% ============================================================
%  ENCODING E FONT
% ============================================================
\usepackage{fontspec}
\setmainfont{TeX Gyre Pagella}
% ============================================================
%  LINGUE
% ============================================================
\usepackage{polyglossia}
\setmainlanguage{italian}
\setotherlanguage[variant=ancient]{greek}

% ============================================================
%  MATEMATICA
% ============================================================
\usepackage{amsmath}
\usepackage{amssymb}
\usepackage{amsfonts}
\usepackage{amsthm}
\usepackage{mathtools}
\usepackage{physics}
\usepackage{yhmath}
\usepackage{cancel}
\usepackage{bbm}

% unicode-math va caricato dopo amsmath/amssymb/amsfonts (altrimenti conflitti
% tipo "\eth already defined"); fornisce il font matematico per XeLaTeX.
\usepackage{unicode-math}
\setmathfont{TeX Gyre Pagella Math}

%dimensionimatematica
\makeatletter
\DeclareMathSizes{\@xpt}{10}{7.5}{6}    % testo 10pt
\DeclareMathSizes{\@xipt}{10}{7.5}{6}   % testo 11pt
\DeclareMathSizes{\@xiipt}{10}{7.5}{6}  % testo 12pt
\makeatother

% Spaziatura display math
\setlength{\abovedisplayskip}{4pt}
\setlength{\belowdisplayskip}{4pt}
\setlength{\abovedisplayshortskip}{2pt}
\setlength{\belowdisplayshortskip}{2pt}

% Comandi matematici
\newcommand{\R}{\mathbb{R}}
\newcommand{\sse}{ \ \textup{sse} \ }
\newcommand{\tc}{ \textup{ t.c. }}
\newcommand{\under}[1]{_{\textup{#1}}}
\newcommand{\sopra}[1]{^{\textup{#1}}}
\newcommand{\tonde}[1]{\left(#1\right)}
\newcommand{\quadre}[1]{\left[#1\right]}
\newcommand{\graffe}[1]{\left\lbrace #1 \right\rbrace}
\newcommand{\et}{\wedge}
\newcommand{\vel}{\vee}
\newcommand{\aut}{\dot{\vee}}
\newcommand{\quindi}{\quad \Rightarrow \quad}
\newcommand{\ubar}[1]{\underline{\mathord{#1}}}

% ============================================================
%  UNITÀ DI MISURA
% ============================================================
\usepackage{siunitx}
\sisetup{
  detect-all,
  output-decimal-marker = {,},
  per-mode = symbol
}
\DeclareSIPrefix\micro{μ}{-6}
% ============================================================
%  LAYOUT E GEOMETRIA
% ============================================================
\usepackage[centering, includeheadfoot]{geometry}

\newlength{\margineEsterno}
\newlength{\margineInterno}
\newlength{\margineAlto}
\newlength{\margineBasso}

\setlength{\margineEsterno}{3.5cm}
\setlength{\margineInterno}{1cm}
\setlength{\margineAlto}{1cm}
\setlength{\margineBasso}{1cm}

\geometry{
  paperwidth  = 170mm,
  paperheight = 240mm,
  top         = \margineAlto,
  bottom      = \margineBasso,
  inner       = \margineInterno,
  outer       = \margineEsterno,
  twoside,
}

% Dimensioni margine per QR e note
\newlength{\maxQR}
\newlength{\sizeQR}
\newlength{\largQR}

\setlength{\maxQR}{2.5cm}
\setlength{\marginparsep}{3mm}
\setlength{\marginparwidth}{\dimexpr\margineEsterno - \marginparsep\relax}
\setlength{\largQR}{\dimexpr\marginparwidth - 6mm\relax}

\newcommand{\calcsizeQR}{%
  \setlength{\sizeQR}{\dimexpr\marginparwidth - 10pt\relax}%
  \ifdim\sizeQR>\maxQR \setlength{\sizeQR}{\maxQR}\fi
}

% Spaziatura testo
\usepackage{setspace}
\setstretch{1.18}
\usepackage{parskip}
\setlength{\parindent}{0pt}
\setlength{\parskip}{0.25\baselineskip plus 1pt minus 1pt}

%Sforare le immagini nel margine esterno
\newcommand{\sforaEsterno}[1]{%
  \checkoddpage
  \ifoddpage
    \begin{adjustwidth}{0pt}{-\margineEsterno}#1\end{adjustwidth}%
  \else
    \begin{adjustwidth}{-\margineEsterno}{0pt}#1\end{adjustwidth}%
  \fi
}

\raggedbottom

% ============================================================
%  INTESTAZIONI E PIÈ DI PAGINA
% ============================================================
\usepackage{fancyhdr}
\pagestyle{fancy}
\fancyhf{}
\fancyfoot[C]{\thepage}
\renewcommand{\headrulewidth}{0pt}
\renewcommand{\footrulewidth}{0pt}
\let\cleardoublepage\clearpage
\renewcommand{\cleardoublepage}{\clearpage}

% ============================================================
%  TITOLI E SEZIONI
% ============================================================
\usepackage{titlesec}
\usepackage{tocloft}

% Part
\renewcommand{\part}[1]{%
  \refstepcounter{part}
  \clearpage
  \thispagestyle{empty}
  \begin{center}
    \Huge\bfseries Part \thepart \\[1em] #1
  \end{center}
  \vspace{2em}
  \addcontentsline{toc}{part}{Part \thepart: #1}
}

% Helper: salva il nome del capitolo E lo stampa
\newcommand{\savechaptertitle}[1]{\gdef\currentchaptername{#1}#1}

\titleformat{\chapter}[block]
  {\normalfont\LARGE\bfseries\color{cyan}}
  {\rule{\textwidth}{.5pt}\\[0.1em]\centering\Roman{chapter}.\,}
  {1em}
  {\huge\scshape\centering\savechaptertitle}
  [\vspace{-0.2em}\rule{\textwidth}{0.5pt}]

% ============================================================
%  CAPITOLO CON IMMAGINE OPZIONALE
% ============================================================
\newcommand{\chapterimage}[3]{%
  \clearpage
  \thispagestyle{fancy}%
  \begin{adjustwidth}{-\margineEsterno}{0pt}%
    \includegraphics[width=\dimexpr\textwidth+\margineEsterno\relax]{#1}%
    \par\vspace{2pt}%
    {\small\itshape\color{gray} #2}%
  \end{adjustwidth}%
  \vspace{0.5em}%
  \refstepcounter{chapter}%
  \gdef\currentchaptername{#3}%
  \addcontentsline{toc}{chapter}{\protect\numberline{\thechapter}#3}%
  \begin{center}
    {\color{cyan}\rule{\textwidth}{.5pt}}\\[0.2em]
    {\color{cyan}\LARGE\bfseries\Roman{chapter}.\,}\\[0.1em]
    {\color{cyan}\huge\scshape #3}\\[-0.3em]
    {\color{cyan}\rule{\textwidth}{.5pt}}
  \end{center}
  \vspace{0.5em}%
}

\titleformat{\section}[block]
  {\normalfont\Large\bfseries\color{cyan}}{\thesection}{1em}{}

\titleformat{\subsection}[block]
  {\normalfont\large\bfseries\color{cyan}}{\thesubsection}{1em}{}

\titleformat{\subsubsection}[block]
  {\normalfont\bfseries\color{cyan}}{\thesubsubsection}{1em}{}

\titleformat{\paragraph}[block]
  {\normalfont\bfseries\color{cyan}}{}{0em}{}

\titlespacing*{\chapter}   {0pt}{-30pt}{8pt}
\titlespacing*{\section}   {0pt}{0.8\baselineskip}{0.3\baselineskip}
\titlespacing*{\subsection}{0pt}{0.6\baselineskip}{0.2\baselineskip}
\setlength{\cftbeforechapskip}{4pt}

% ============================================================
%  LISTE
% ============================================================
\usepackage[inline]{enumitem}
\setlist[itemize,1]{label=\textcolor{gray!70}{\rule{0.8ex}{0.8ex}}}
\setlist{itemsep=0pt, parsep=0pt, topsep=3pt, partopsep=0pt}
\setlength{\topsep}{2pt}
\newlist{myenum}{enumerate*}{1}
\setlist[myenum]{label=\arabic*, itemjoin={\hspace{1.5cm}}}
\newlist{sceltamultipla}{enumerate}{1}
\setlist[sceltamultipla,1]{%
  label=\textcolor{gray}{\fbox{\textbf{\footnotesize\makebox[1.2em]{\Alph*}}}},
  leftmargin=*,
  itemsep=1pt,
  topsep=2pt
}

% ============================================================
%  COLORI
% ============================================================
\usepackage{xcolor}

\definecolor{ciano}           {RGB}{0,   180, 200}
\definecolor{sfondoImportante}{RGB}{255, 252, 220}
\definecolor{coloreAnimazione}{RGB}{0,   150, 136}
\definecolor{sfondoAnimazione}{RGB}{225, 245, 243}
\definecolor{gold}            {RGB}{255, 215,   0}
\definecolor{copper}          {RGB}{104, 115,  51}

\definecolor{buildSub}  {HTML}{FFF0F5}  \definecolor{buildTitle} {HTML}{C71585}
\definecolor{readSub}   {HTML}{F2F2F2}  \definecolor{readTitle}  {HTML}{333333}
\definecolor{deriveSub} {HTML}{FFF9E6}  \definecolor{deriveTitle}{HTML}{B8860B}
\definecolor{exSub}     {HTML}{F0F7FF}  \definecolor{exTitle}    {HTML}{005C99}
\definecolor{impSub}    {HTML}{FFF0F0}  \definecolor{impTitle}   {HTML}{CC0000}
\definecolor{warnSub}   {HTML}{FFF0F0}  \definecolor{warnTitle}  {HTML}{B22222}

\definecolor{sfondoRiassunto}{RGB}{215, 242, 240}
\definecolor{intestRiassunto}{RGB}{  0, 145, 132}
\definecolor{rsezColor}      {RGB}{  0, 160, 146}

% ============================================================
%  TIKZ E GRAFICI
% ============================================================
\usepackage{tikz}
\usepackage{tikz-3dplot}
\usepackage{pgfplots}
\usepackage{circuitikz}
\usepackage{tikzsymbols, fontawesome5}

\usetikzlibrary{
  decorations.pathreplacing, decorations.markings, decorations.pathmorphing,
  angles, quotes, arrows.meta, calc, positioning,
  shapes, shapes.symbols, shadows, fit,
  backgrounds, matrix, patterns, tikzmark
}

\pgfplotsset{compat=1.17}
\usepgfplotslibrary{colormaps}
\tdplotsetmaincoords{70}{110}

% Simbolo caution
\newcommand{\caution}[1]{%
  \begingroup
  \pgfmathsetmacro{\s}{#1}%
  \pgfmathsetmacro{\Rr}{1.6*\s}%
  \pgfmathsetmacro{\Lwd}{2*\s}%
  \pgfmathsetmacro{\FSpt}{40*\s}%
  \pgfmathsetmacro{\BLSpt}{1.2*\FSpt}%
  \tikz[baseline=(bang.base)]{
    \draw[line width=\Lwd mm, red, fill=yellow!30]
      (90:\Rr) -- (210:\Rr) -- (330:\Rr) -- cycle;
    \node (bang) at (0,0)
      {\fontsize{\FSpt pt}{\BLSpt pt}\selectfont\bfseries !};
  }%
  \endgroup
}

% ============================================================
%  BOX: TCOLORBOX
% ============================================================
\usepackage[many]{tcolorbox}

\tcbset{
  mystyle/.style={
    enhanced, breakable,
    boxrule=0pt, frame hidden, arc=2mm,
    top=4mm, bottom=4mm, left=4mm, right=4mm,
    fonttitle=\bfseries\sffamily, coltitle=white,
    attach boxed title to top left={yshift=-2mm, xshift=0mm},
    boxed title style={
      empty,
      top=1.5mm, bottom=1.5mm, left=3mm, right=8mm,
      underlay={
        \fill[#1, rounded corners=2mm]
          (frame.south west) -- (frame.north west) --
          ([xshift=-5mm]frame.north east) -- (frame.south east) -- cycle;
      }
    }
  }
}

\newtcolorbox{mdwarning}[1][]{
  mystyle={warnTitle}, colback=warnSub,
  title={\ifstrempty{#1}{Attenzione}{#1}}
}
\newtcolorbox{mdbuild}[1][]{
  mystyle={buildTitle}, colback=buildSub,
  title={\ifstrempty{#1}{Costruire la formula}{#1}}
}
\newtcolorbox{mdread}[1][]{
  mystyle={readTitle}, colback=readSub,
  title={\ifstrempty{#1}{Leggere la formula}{#1}}
}
\newtcolorbox{mdderive}[1][]{
  mystyle={deriveTitle}, colback=deriveSub,
  title={\ifstrempty{#1}{Derivare la formula}{#1}}
}
\newtcolorbox{mdexample}[1][]{
  mystyle={exTitle}, colback=exSub,
  title={\ifstrempty{#1}{Esempio guidato}{#1}}
}

\newtcolorbox{mdimportant}[2][]{%
  enhanced, breakable,
  boxrule=0pt, frame hidden,
  colback=sfondoImportante, arc=1mm,
  top=4mm, bottom=4mm, left=6mm, right=4mm,
  overlay={
    \checkoddpage
    \ifoddpage
      \draw[ciano, line width=2pt] (frame.south east)--(frame.north east);
      \draw[ciano, line width=2pt] (frame.east)--([xshift=5mm]frame.east);
      \draw[ciano, fill=white, line width=2pt] ([xshift=5mm]frame.east) circle (4pt);
      \node[anchor=west, text=ciano, font=\bfseries\sffamily\small,
            align=left, text width=\dimexpr\largQR-2mm\relax]
        at ([xshift=7mm]frame.east) {#2};
    \else
      \draw[ciano, line width=2pt] (frame.south west)--(frame.north west);
      \draw[ciano, line width=2pt] ([xshift=-5mm]frame.west)--(frame.west);
      \draw[ciano, fill=white, line width=2pt] ([xshift=-5mm]frame.west) circle (4pt);
      \node[anchor=east, text=ciano, font=\bfseries\sffamily\small,
            align=right, text width=\dimexpr\largQR-2mm\relax]
        at ([xshift=-7mm]frame.west) {#2};
    \fi
  },
  #1
}

\newsavebox{\mdanimationQRbox}

\newtcolorbox{mdanimation}[2][Animazione]{%
  enhanced, breakable,
  boxrule=0pt, frame hidden,
  colback=sfondoAnimazione, arc=2mm,
  top=4mm, bottom=4mm, left=4mm, right=4mm,
  fonttitle=\bfseries\sffamily,
  coltitle=white,
  title={#1},
  attach boxed title to top left={yshift=-2mm, xshift=0mm},
  boxed title style={
    empty,
    top=1.5mm, bottom=1.5mm, left=3mm, right=8mm,
    underlay={
      \fill[coloreAnimazione, rounded corners=2mm]
        (frame.south west) -- (frame.north west) --
        ([xshift=-5mm]frame.north east) -- (frame.south east) -- cycle;
    }
  },
  overlay={
    \ifx\relax#2\relax\else
      \checkoddpage
      \calcsizeQR
      \ifoddpage
        \fill[sfondoAnimazione]
          ([xshift=-1mm, yshift=-2mm]frame.north east) --
          ([xshift=3mm, yshift=1mm]frame.north east) --
          ([xshift=3mm, yshift=-\the\sizeQR-0pt]frame.north east) --
          ([xshift=-1mm, yshift=-\the\sizeQR+3mm]frame.north east) -- cycle;
        \node[anchor=north west, draw=coloreAnimazione, line width=1.5pt,
              fill=sfondoAnimazione, inner sep=4pt, rounded corners=1mm]
          at ([xshift=3mm, yshift=2mm]frame.north east) {%
            \begin{minipage}{\the\largQR}\centering
              \qrcode[height=\the\sizeQR]{#2}
            \end{minipage}%
          };
      \else
        \fill[sfondoAnimazione]
          ([xshift=1mm, yshift=-2mm]frame.north west) --
          ([xshift=-3mm, yshift=1mm]frame.north west) --
          ([xshift=-3mm, yshift=-\the\sizeQR-0pt]frame.north west) --
          ([xshift=1mm, yshift=-\the\sizeQR+3mm]frame.north west) -- cycle;
        \node[anchor=north east, draw=coloreAnimazione, line width=1.5pt,
              fill=sfondoAnimazione, inner sep=4pt, rounded corners=1mm]
          at ([xshift=-3mm, yshift=2mm]frame.north west) {%
            \begin{minipage}{\the\largQR}\centering
              \qrcode[height=\the\sizeQR]{#2}
            \end{minipage}%
          };
      \fi
    \fi
  },
}

% ============================================================
%  BOX: MDFRAMED
% ============================================================
\usepackage{mdframed}

\newenvironment{cyanmargin}[1][]
  {\begin{mdframed}[topline=false,bottomline=false,rightline=false,
    linecolor=cyan,linewidth=.05cm,frametitlefont=\itshape]#1\ignorespaces}
  {\end{mdframed}}

\newenvironment{mdexercise}[1][]
  {\begin{mdframed}[topline=false,bottomline=false,rightline=false,
    linecolor=cyan,linewidth=.05cm,frametitlefont=\itshape]
   \textbf{Esercizio. }#1\ignorespaces}
  {\end{mdframed}}

\newenvironment{mdexam}[1][]
  {\begin{mdframed}[topline=false,bottomline=false,rightline=false,
    linecolor=magenta,linewidth=2,frametitlefont=\itshape]
   \textbf{Allenamento per la verifica. }#1\ignorespaces}
  {\end{mdframed}}

\newenvironment{mdproof}[1][]
  {\begin{mdframed}[topline=false,bottomline=false,rightline=false,
    linecolor=black,frametitle=Proof.,frametitlefont=\itshape]#1\ignorespaces}
  {\unskip\begin{flushright}\qedsymbol\end{flushright}\end{mdframed}}

\newenvironment{nestedproof}[1][\small]
  {\begin{mdframed}[topline=false,bottomline=false,rightline=false,
    linecolor=black,frametitle=Proof.,frametitlefont=\itshape]#1\ignorespaces}
  {\unskip\begin{flushright}\qedsymbol\end{flushright}\end{mdframed}}

% ============================================================
%  TEOREMI
% ============================================================
\newtheoremstyle{flemma}
  {3pt}{3pt}
  {\itshape\fontfamily{qpl}\selectfont}
  {1pt}{\bfseries}{ \space}{ }{}

\theoremstyle{definition}
\newtheorem*{definition}  {Definizione}
\newtheorem*{proposition} {Proposizione}
\newtheorem*{example}     {Esempio}
\newtheorem*{remark}      {Osservazione}
\newtheorem*{nota}        {Nota}
\newtheorem*{application} {Applicazione}

\theoremstyle{flemma}
\newtheorem*{property}  {Proprietà}
\newtheorem*{law}       {Legge}
\newtheorem*{lemma}     {Lemma}
\newtheorem*{theorem}   {Teorema}
\newtheorem*{corollary} {Corollario}
\newtheorem*{principle} {Principio}

\renewcommand\qedsymbol{\fontfamily{qpl}\selectfont\emph{Q.E.D.}}

% ============================================================
%  ESERCIZI E SOLUZIONI (ambienti classici)
% ============================================================
\usepackage{needspace}
\usepackage{xparse}
\usepackage{pifont}
\usepackage{environ}

\newcounter{exercise}
\renewcommand{\theexercise}{\arabic{exercise}}

% Stelline — nere, senza fbox
\newcommand{\onestar}  {\ding{72}\ding{73}\ding{73}\space}
\newcommand{\twostars} {\ding{72}\ding{72}\ding{73}\space}
\newcommand{\threestars}{\ding{72}\ding{72}\ding{72}\space}

\newcommand{\setexdifficulty}[1]{%
  \IfNoValueTF{#1}
    {\def\exdifficulty{\twostars}}
    {\def\exdifficulty{#1}}%
}

\NewDocumentEnvironment{exercise}{ o }{%
  \refstepcounter{exercise}\par\medskip
  \setexdifficulty{#1}\noindent
  \hypertarget{ex:\theexercise}{}\label{ex:\theexercise}%
  \textbf{Esercizio \theexercise\ }\exdifficulty\quad
  \hyperlink{sol:\theexercise}{\textit{(vai alla soluzione)}}\quad\ignorespaces
}{\par\medskip}

\NewDocumentEnvironment{question}{ o }{%
  \refstepcounter{exercise}\par\medskip
  \setexdifficulty{#1}\noindent
  \hypertarget{ex:\theexercise}{}\label{ex:\theexercise}%
  \textbf{Domanda di Teoria \theexercise\ }\exdifficulty\quad
  \hyperlink{sol:\theexercise}{\textit{(vai alla soluzione)}}\quad\ignorespaces
}{\par\medskip}

\NewDocumentEnvironment{exercisenosol}{ o }{%
  \refstepcounter{exercise}\par\medskip
  \setexdifficulty{#1}\noindent
  \hypertarget{ex:\theexercise}{}\label{ex:\theexercise}%
  \textbf{Esercizio \theexercise\ }\exdifficulty\quad\ignorespaces
}{\par\medskip}

\NewDocumentEnvironment{questionnosol}{ o }{%
  \refstepcounter{exercise}\par\medskip
  \setexdifficulty{#1}\noindent
  \hypertarget{ex:\theexercise}{}\label{ex:\theexercise}%
  \textbf{Domanda di teoria \theexercise\ }\exdifficulty\quad\ignorespaces
}{\par\medskip}

\NewDocumentEnvironment{expsolution}{ o }{%
  \Needspace{6\baselineskip}%
  \IfNoValueTF{#1}{\edef\resnum{\theexercise}}{\edef\resnum{#1}}%
  \par\bigskip
  \hypertarget{sol:\resnum}{}\label{sol:\resnum}%
  \noindent\textbf{Soluzione \hyperlink{ex:\resnum}{\resnum}}\quad
  \hyperlink{ex:\resnum}{\textit{(torna all'esercizio)}}\par\medskip
}{\par\medskip}

% solution: comportamento duale (inline grigio fuori da rigaesercizi,
%           accumulatore dentro rigaesercizi — vedi sotto)
\newenvironment{solution}
  {\par\noindent\hfill\small\color{gray}}
  {\par}

% ============================================================
%  MODALITÀ ESERCIZI (due colonne + testata ciano)
% ============================================================
\usepackage{multicol}
\usepackage{colortbl}

\newcommand{\exheadrule}[2]{%
  \setlength{\headheight}{20pt}%
  \fancyhead[RE]{%
    \colorbox{ciano}{%
      \parbox[c][14pt][c]{\textwidth}{%
        \color{white}\bfseries\sffamily\small #1\hfill}}}%
  \fancyhead[LO]{%
    \colorbox{ciano}{%
      \parbox[c][14pt][c]{\textwidth}{%
        \color{white}\bfseries\sffamily\small\hfill #2}}}%
  \fancyhead[RO,LE]{}%
  \renewcommand{\headrulewidth}{0pt}%
}

\newif\ifinsidemodalita
\insidemodalitafalse

\newenvironment{modalitaesercizi}{%
  \insidemodalitatrue
  \edef\exChapNum{\thechapter}%
  \edef\exChapName{\currentchaptername}%
  \renewcommand{\Needspace}[1]{}%
  \newgeometry{
    top=0cm, bottom=.8cm,
    inner=1.2cm, outer=1.2cm,
    includeheadfoot,
  }%
  \pagestyle{fancy}\fancyhf{}\fancyfoot[C]{\thepage}%
  \exheadrule%
    {CAPITOLO \exChapNum \enspace\textbar\enspace \exChapName}%
    {CAPITOLO \exChapNum \enspace\textbar\enspace ESERCIZI}%
  \setlength{\columnsep}{1cm}%
  \raggedcolumns
  \begin{multicols}{2}%
}{%
  \end{multicols}%
  \restoregeometry%
  \pagestyle{fancy}\fancyhf{}\fancyhead{}\fancyfoot[C]{\thepage}%
  \renewcommand{\headrulewidth}{0pt}%
  \insidemodalitafalse
}

\newcommand{\pausamulticol}{%
  \ifinsidemodalita\end{multicols}\par\vspace{2pt}\fi
}
\newcommand{\riprendimulticol}{%
  \ifinsidemodalita\par\vspace{2pt}\begin{multicols}{2}\fi
}

% ============================================================
%  AMBIENTE truthorfalse
%  Da aggiungere nel preambolo, dopo il caricamento di
%  tcolorbox (già presente) e enumitem (già presente).
%
%  Uso:
%    \begin{truthorfalse}
%      \tfitem{testo dell'affermazione}
%    \end{truthorfalse}
%
%  Le caselle V / F vengono inserite automaticamente
%  a destra di ogni voce.
% ============================================================

% Casellina grigia piccola (V oppure F).
% Si usa \newtcbox invece di \tcbox[...]{...} inline:
% \tcbox con opzioni esplicite dentro \newcommand causa
% conflitti nel parsing degli argomenti di tcolorbox.
% \newtcbox pre-registra le opzioni ed è privo di ambiguità.
\newtcbox{\tfvf}{%
  on line,
  arc=0.8pt, outer arc=0.8pt,
  boxrule=0.3pt, boxsep=0pt,
  left=2pt, right=2pt, top=0.8pt, bottom=0.8pt,
  colframe=gray!40, colback=gray!10,
  fontupper={\fontsize{6}{6}\selectfont\bfseries\color{gray!65}}%
}

% Singola voce: testo → caselle V F allineate a destra
\newcommand{\tfitem}[1]{%
  \item #1\nobreak\hfill\tfvf{V}\kern2pt\tfvf{F}%
}

% Lista vera/falsa
\newenvironment{truthorfalse}{%
  \begin{enumerate}[%
    label  = \textbf{\alph*.},%
    leftmargin = *,%
    itemsep = 3pt,%
    topsep  = 2pt,%
    parsep  = 0pt%
  ]%
}{%
  \end{enumerate}%
}
% ============================================================
%  Flag per raccolta soluzioni
% ============================================================
% ============================================================
%  *** SOSTITUIRE TUTTO da qui fino a fine \newenvironment{exrcell} ***
% ============================================================

% Flag per raccolta soluzioni
\newif\ifincollectmode  \incollectmodefalse
\newif\iffirstsolution  \firstsolutiontrue

% ============================================================
%  CORREZIONE BUG soluzioni duplicate con rigaesercizi{3}
%
%  Causa: il vecchio meccanismo con \newtoks memorizzava
%  \the\exrsoltemp come riferimento LAZY dentro \exrsoltoks.
%  Quando la terza soluzione sovrascriveva \exrsoltemp,
%  entrambi i riferimenti espandevano lo stesso valore:
%  risultato [sol1, sol3, sol3] invece di [sol1, sol2, sol3].
%
%  Fix: ogni soluzione viene copiata permanentemente in
%  \exrsol@1, \exrsol@2, \exrsol@3  via \let, che fotografa
%  \BODY al momento della chiamata.  Ridefinizioni successive
%  di \BODY non toccano le copie già salvate.
% ============================================================

\newcounter{exrsolcount}   % conta le soluzioni della riga corrente

\RenewEnviron{solution}{%
  \ifincollectmode
    \iffirstsolution
      \global\firstsolutionfalse
    \fi
    \stepcounter{exrsolcount}%
    % Copia permanente: \exrsol@N ottiene il significato corrente
    % di \BODY; eventuali ridefinizioni future di \BODY non la toccano
    \expandafter\global\expandafter\let
      \csname exrsol@\theexrsolcount\endcsname\BODY
  \else
    \par\noindent\hfill\small\color{gray}\BODY\par
  \fi
}

% ============================================================
%  Dimensioni e box ausiliari (invariati)
% ============================================================
\newsavebox{\exrfulltmpbox}
\newsavebox{\exrsolbox}
\newsavebox{\exrstarbox}
\newlength{\exrlabelwidth}
\newlength{\exrcellwidth}
\newlength{\exrsolcolwidth} \setlength{\exrsolcolwidth}{3.5cm}
\newlength{\exrcellsep}     \setlength{\exrcellsep}{6mm}
\newlength{\exrlabelgap}    \setlength{\exrlabelgap}{4mm}
\newlength{\exrtotalgap}
\newcounter{exrowcell}
\newcounter{exrowmaxcolsctr}

% ============================================================
%  Stampa le soluzioni accumulate come  [sol1; sol2; sol3]
%  Gestisce fino a 3 soluzioni per riga (massimo supportato).
% ============================================================
\newcommand{\exrprintsol}{%
  \savebox{\exrsolbox}{%
    \color{gray!70!black}%
    \ifnum\value{exrsolcount}>0\relax
      \csname exrsol@1\endcsname
      \ifnum\value{exrsolcount}>1\relax
        ;\space\csname exrsol@2\endcsname
        \ifnum\value{exrsolcount}>2\relax
          ;\space\csname exrsol@3\endcsname
        \fi
      \fi
    \fi
  }%
  \color{gray!70!black}%
  $\left[\rule[-\dp\exrsolbox]{0pt}%
    {\dimexpr\ht\exrsolbox+\dp\exrsolbox\relax}\right.$%
  \usebox{\exrsolbox}%
  $\left.\rule[-\dp\exrsolbox]{0pt}%
    {\dimexpr\ht\exrsolbox+\dp\exrsolbox\relax}\right]$%
}

\newenvironment{rigaesercizi}[2][\twostars]{%
  \setcounter{exrowcell}{0}%
  \setcounter{exrowmaxcolsctr}{#2}%
  \setcounter{exrsolcount}{0}%
  \global\firstsolutiontrue%
  \incollectmodetrue%
  \refstepcounter{exercise}%
  \setexdifficulty{#1}%
  \hypertarget{ex:\theexercise}{}\label{ex:\theexercise}%
  \sbox{\exrstarbox}{{\scriptsize\exdifficulty}}%
  \setlength{\exrlabelwidth}{\dimexpr\wd\exrstarbox + 6pt\relax}%
  \setlength{\exrtotalgap}{#2\exrcellsep}%
  \advance\exrtotalgap by -\exrcellsep%
  \ifnum#2=1\relax
    \setlength{\exrcellwidth}{\dimexpr
      \linewidth - \exrlabelwidth - \exrlabelgap\relax}%
  \else
    \setlength{\exrcellwidth}{\dimexpr
      ( \linewidth
        - \exrlabelwidth
        - \exrlabelgap
        - \exrsolcolwidth
        - \exrcellsep
        - \exrtotalgap
      ) / #2\relax}%
  \fi
  \par\nopagebreak\vspace{2pt}%
  % \vbox\bgroup\noindent\leavevmode
  \vbox\bgroup\noindent\leavevmode\parfillskip=0pt
  {\setlength{\fboxsep}{0pt}%
   \vtop{\offinterlineskip
     \hbox{\colorbox{intestRiassunto}{%
       \color{white}\bfseries\small
       \makebox[\exrlabelwidth][c]{\rule{0pt}{11pt}\theexercise}}}%
     \hbox{\makebox[\exrlabelwidth][c]{\scriptsize\exdifficulty}}%
   }}%
  \hspace{\exrlabelgap}\small\raggedright\parfillskip=0pt
}{%
  \iffirstsolution\else
    \unskip\nobreak\hfill\allowbreak\null\hfill\mbox{\exrprintsol}%
  \fi
  \par\egroup
  \vspace{4pt}%
  \incollectmodefalse%
}

\newenvironment{exrcell}{%
  \stepcounter{exrowcell}%
  \ifnum\value{exrowmaxcolsctr}=1\relax
    \ignorespaces
  \else
    \vtop\bgroup
      \hsize=\exrcellwidth
      \binoppenalty=10000
      \relpenalty=10000
      \small\raggedright\ignorespaces
  \fi
}{%
  \ifnum\value{exrowmaxcolsctr}=1\relax
    \unskip
  \else
    \egroup
    \ifnum\value{exrowcell}<\value{exrowmaxcolsctr}\relax
      \hspace{\exrcellsep}%
    \fi
  \fi
}

% ============================================================
%  *** FINE sezione da sostituire ***
% ============================================================
% ============================================================
%  SEZIONI INTERNE AL RIASSUNTO
% ============================================================
\newcommand{\rsection}[1]{%
  \par\vspace{0.55\baselineskip}%
  \begin{tcolorbox}[
    enhanced, breakable,
    colback=rsezColor, colframe=rsezColor,
    arc=2mm, boxrule=0pt,
    top=2mm, bottom=2mm, left=5mm, right=5mm,
  ]
    \color{white}\bfseries\sffamily\large #1%
  \end{tcolorbox}%
  \vspace{0.1\baselineskip}%
}

\newcommand{\rsubsection}[1]{%
  \par\vspace{0.35\baselineskip}%
  \noindent{\color{intestRiassunto}\bfseries\sffamily #1}%
}

\newcommand{\rkeyword}[1]{%
  \textcolor{intestRiassunto}{\bfseries #1}%
}

% ============================================================
%  AMBIENTE RIASSUNTO
% ============================================================
\newenvironment{riassunto}{%
  \clearpage
  \pagecolor{sfondoRiassunto}%
  \newgeometry{
    top    = 1.8cm,
    bottom = 1.8cm,
    inner  = 1.5cm,
    outer  = 1.5cm,
    includeheadfoot,
  }%
  \pagestyle{fancy}\fancyhf{}%
  \setlength{\headheight}{22pt}%
  \fancyhead[LE,RO]{%
    \colorbox{intestRiassunto}{%
      \parbox[c][16pt][c]{\textwidth}{%
        \color{white}\bfseries\sffamily\small
        \hspace{3mm}CAPITOLO~\thechapter
        \enspace\textbar\enspace
        \MakeUppercase{\currentchaptername}%
        \hfill RIASSUNTO\hspace{3mm}%
      }%
    }%
  }%
  \fancyhead[LO,RE]{}%
  \renewcommand{\headrulewidth}{0pt}%
  \fancyfoot[C]{\color{intestRiassunto}\thepage}%
  \par\vspace{0.6em}%
  \begin{center}
    {\color{intestRiassunto!50!black}\rule{0.18\textwidth}{1pt}}%
    \quad
    {\color{intestRiassunto}\fontsize{27}{31}\selectfont\bfseries\scshape Riassunto}%
    \quad
    {\color{intestRiassunto!50!black}\rule{0.18\textwidth}{1pt}}\\[0.35em]
    {\color{rsezColor!70!white}\fontsize{9}{11}\selectfont\sffamily
      Capitolo~\thechapter\ \textbullet\ \currentchaptername}%
  \end{center}%
  \vspace{0.7em}%
}{%
  \clearpage
  \pagecolor{white}%
  \restoregeometry%
  \pagestyle{fancy}%
  \fancyhf{}\fancyhead{}\fancyfoot[C]{\thepage}%
  \renewcommand{\headrulewidth}{0pt}%
}

% ============================================================
%  AMBIENTE VERIFICA SUBITO
% ============================================================
\newenvironment{verificasubito}[2][Verifica subito!]{%
  \def\vsncols{#2}%
  \renewcommand{\Needspace}[1]{}%
  \checkoddpage%
  \ifoddpage%
    \def\sforoSin{0cm}%
    \def\sforoDes{0cm}%
  \else%
    \def\sforoSin{0cm}%
    \def\sforoDes{0cm}%
  \fi%
  \begin{tcolorbox}[
    enhanced, breakable,
    colback=white, colframe=ciano,
    boxrule=1.5pt, arc=3mm,
    top=4mm, bottom=4mm, left=4mm, right=4mm,
    enlarge left by=\sforoSin,
    enlarge right by=\sforoDes,
    title={\large\bfseries\sffamily #1},
    coltitle=white,
    attach boxed title to top left={yshift=-2mm, xshift=4mm},
    boxed title style={
      empty,
      top=1.5mm, bottom=1.5mm, left=3mm, right=8mm,
      underlay={
        \fill[ciano, rounded corners=2mm]
          (frame.south west) -- (frame.north west) --
          ([xshift=-5mm]frame.north east) -- (frame.south east) -- cycle;
      }
    },
  ]%
  \setlength{\columnsep}{1cm}%
  \raggedcolumns%
  \ifnum\vsncols>1\relax
    \begin{multicols}{\vsncols}%
  \fi
}{%
  \ifnum\vsncols>1\relax
    \end{multicols}%
  \fi
  \end{tcolorbox}%
}

% ============================================================
%  QR NEL MARGINE
% ============================================================
\usepackage{qrcode}
\usepackage{marginnote}
\usepackage{marginfix}

\newcommand{\marginqr}[2][]{%
  \par\Needspace{6\baselineskip}\calcsizeQR
  \marginpar{\centering
    \begin{tcolorbox}[enhanced, boxrule=1.5pt, colframe=ciano, colback=white,
      arc=1mm, top=4pt, bottom=4pt, left=3pt, right=3pt,
      boxsep=0pt, width=\marginparwidth]
      \if\relax\detokenize{#1}\relax\else
        {\footnotesize\setlength{\parskip}{0pt}\setstretch{1.0}%
         \raggedright\sffamily #1\par\vspace{3pt}}
      \fi
      \centering\qrcode[height=\sizeQR]{#2}%
    \end{tcolorbox}}%
}


% ============================================================
%  AMBIENTE PROVATU — box nel margine, colori riassunto
% ============================================================
\newsavebox{\provatubox}
\newenvironment{provatu}[1][Prova tu!]{%
  \begin{lrbox}{\provatubox}%
    \begin{minipage}{\marginparwidth}%
      \begin{mdframed}[
        skipabove=0pt, skipbelow=0pt,
        linecolor=intestRiassunto,
        linewidth=1.2pt,
        backgroundcolor=sfondoRiassunto,
        frametitle={\color{white}\small\bfseries\sffamily\hspace{2pt}#1},
        roundcorner=4pt,
        frametitlebackgroundcolor=intestRiassunto,
        frametitlerule=false,
        frametitleaboveskip=2pt,
        frametitlebelowskip=2pt,
        innertopmargin=4pt,
        innerbottommargin=4pt,
        innerleftmargin=4pt,
        innerrightmargin=4pt,
      ]%
      \small\raggedright\ignorespaces
}{%
      \end{mdframed}%
    \end{minipage}%
  \end{lrbox}%
  \marginnote{\usebox{\provatubox}}%
  \ignorespaces
}{\ignorespacesafterend}

% ============================================================
%  TESTO E TIPOGRAFIA
% ============================================================
\usepackage{csquotes}
\usepackage[normalem]{ulem}
\usepackage{soul}
\usepackage{tipa}

\setulcolor{gray}
\renewcommand{\CancelColor}{\color{red}}

\newcommand{\schwa}{ə}
\newcommand{\Schwa}{Ə}
\newcommand{\lette}[1]{\enquote*{#1}}
\newcommand{\virgo}[1]{\enquote*{#1}}
\newcommand{\definizione}[1]{\textcolor{red}{\textbf{#1}}}
\newcommand{\citaz}[2]{\guillemotleft#1\guillemotright\begin{flushright}#2\end{flushright}}
\newcommand{\cit}[1]{\guillemotleft#1\guillemotright}
\newcommand{\boh}[1]{\textcolor{red}{#1}}
\newcommand{\bianco}[1]{\textcolor{white}{#1}}

% ============================================================
%  PACCHETTI GRAFICI E VARI
% ============================================================
\usepackage{graphicx}
\usepackage{array}
\usepackage{booktabs}
\usepackage{placeins}
\usepackage{titling}
\usepackage{float}
\usepackage{wrapfig}
\usepackage{sidecap}
\usepackage{subfig}
\usepackage{subcaption}
\usepackage{floatflt}
\usepackage{animate}
\usepackage{standalone}
\usepackage{listings}
\usepackage{footmisc}
\usepackage{xstring}
\usepackage{ifthen}
\usepackage[strict]{changepage}
\usepackage[margin=1cm]{caption}

% ============================================================
%  HYPERREF (sempre per ultimo)
% ============================================================
\usepackage[colorlinks]{hyperref}
\hypersetup{linkcolor=magenta}


\begin{document}

\tableofcontents

\chapter{Equazioni di Secondo Grado}
Sto facendo delle lezioni sulle equazioni di Secondo Grado, e utilizzo il libro \textit{Le idee della Matematica (Algebra e Geometria, vol. 3)}. 

\section{Introduzione alle equazioni di secondo grado}

\subsection{Ripasso eq 1 grado}
Chiamare alla lavagna e chiedere di risolvere equazioni di primo grado, generandone sempre di nuove finché non vedo che è chiaro per tutti.
\begin{enumerate}
    \item $-\dfrac 23 x -4 = 2,$
    \item $-\dfrac x2 -2 =2-x$
    \item $-2(x+1)+\dfrac{x-1}2= 1$
\end{enumerate}


\paragraph{Compiti sul ripasso di equazioni di 1 grado}
Improvvisazione

\subsection{Quante soluzioni possono avere le equazioni di 2 grado?}

Cosa succede se prendo una qualunque delle equazioni considerate e metto un \lette{alla seconda} sopra la $x$?
\begin{enumerate}
    \item $-\dfrac 23 x^2 -4 = 2$ $\rightarrow$ nessuna soluzione
    \item $-\dfrac 23 x^2 + 4 = 2$ $\rightarrow$ due soluzioni
    \item $-\dfrac 23 x^2+2=2$ $\rightarrow$ una soluzione.  
\end{enumerate}


Ricordare cosa significa che un certo valore \lette{è una soluzione}.

\begin{mdframed}
    A seconda dei casi, le equazioni di $2^\circ$ grado possono avere 0, 1 o 2 soluzioni.
\end{mdframed}

Benvenuti nell'universo delle equazioni di secondo grado! Cominciate già a capire quello che vi aspetta, eh?

Adesso cerchiamo di definire che cos'è un'equazione di secondo grado.

\section{Le equazioni di secondo grado in forma normale}

\textit{[A voce]} Un'equazione di secondo grado è un'equazione in cui compare una incognita alla seconda e non compare nessuna incognita con una potenza più grande di 2.

\begin{mdframed}
    Un'equazione di $2^\circ$ grado si dice \lette{in forma normale} quando è del tipo
    $$
    ax^2+bx+c = 0, \qquad \text{con }a \ne 0,
    $$
    dove $ax^2$ è il termine quadratico, $bx$ è il termine lineare e\dots
\end{mdframed}

Provare a scrivere alcune equazioni in forma normale e riconoscere i valori di $a$, $b$ e $c$, cercando di prevenire i normali equivoci sull'ordine dei termini tipo $-x+x^2+1=0$, sui segni tipo $-x \rightarrow b=1$, sulle frazioni \lette{sottintese} tipo $\frac{x^2}2 \rightarrow a=\frac 12$, sulle operazioni, etc.

Esempi:
\begin{enumerate}
    \item $2x^2+3-6x=4x^2+1-x$
    \item $-\frac{x^2}2+1=0$
    \item $-2x=3x^2$
    \item $(2y-1)(3y+2)=1$
    \item $(x+1)^2-2x=3x$
    \item $2x^2-3x+4=2x^2$, che non è un'equazione di secondo grado!
\end{enumerate}

Beh, ora che abbiamo imparato a scrivere le equazioni nella forma normale, dovremo pur imparare a risolvere queste equazioni! Ci sono diversi casi...

\paragraph{Compiti sulle equazioni di secondo grado in forma normale}
Se finisco la lezione qui assegno qualche esercizio sul ripasso delle equazioni di primo grado e sulla conversione alla forma normale.



\section{Caso $b=c=0$ (equazioni monomie)}

L'equazione diventa $ax^2=0$. Chiedere se qualcuno ha già idea. Cosa vuol dire che $x=0$ è una soluzione?

Esempi:
\begin{enumerate}
    \item $-\frac{x^2}{2900}=0$
    \item $-3x^2+3=3$
    \item $-2(x-3)^2=0$
    \item $-2(x-3)^2 = 12x-18$
\end{enumerate}

\section{Caso $b=0$ (eq pura)}

L'equazione diventa $ax^2+c=0$. 


Esempi:
\begin{enumerate}
    \item $-\frac23x^2+1=0$
    \item $(2x-1)^2=4-4x+x^2$
    \item $3(2x-1)^2-1=0$
    \item $2x^2+1=0$
    \item $-\frac{x^2}{2}=2$
    \item $(4x+10)^2=-20$
\end{enumerate}

In generale, prendiamo l'equazione
$$
ax^2+c=0
$$
e isoliamo la $x^2$:
$$x^2=-\frac ca \rightarrow x=\pm\sqrt{-\frac ca} .$$

\begin{mdframed}
    Le equazioni pure hanno soluzione solo se $-\frac ca>0$ (cioè se $a$ e $c$ sono discordi). In quel caso, hanno 2 soluzioni, pari a 
    $$
    x_{1,2} = \pm\sqrt{-\frac ca}.
    $$
\end{mdframed}


Proviamo la regola su un esempio! \textit{(Improvvisare)}.

\paragraph{Compiti sulle equazioni pure e monomie}
Se arrivo fin qua posso dare come compiti gli esercizi 22 e 34 sulle equazioni monomie e 21,23,25 su quelle pure.



\section{Caso $c=0$ (equazioni spurie)}

Esempi:
\begin{enumerate}
    \item $2x^2-3x=0$
    \item $3+\frac x4-\frac {x^2}2=3$
    \item $-(x-4)=2(x-4)^2$
\end{enumerate}

Prendi quella in forma generale
$$
ax^2+bx=0
$$
raccogli e ottieni la seguente regola.
\begin{mdframed}
    Per le equazioni senza termine noto (spurie) si raccoglie l'incognita e si utilizza la legge di annullamento del prodotto. Quindi, esistono sempre due soluzioni, di cui una è $0$ e l'altra è $-\frac ba$. 
\end{mdframed}

\paragraph{Compiti sulle equazioni incomplete}
Esercizi 19, 20, 24, 27, 28, 29, 35, 50, 53, 54, 61, 66.

\section{Equazioni di secondo grado complete}

\subsection{La ricetta per la risoluzione di un'equazione di secondo grado}
Bla bla bla\dots

\subsection{Fattorizzazione}

Comincio introducendo alcuni nuovi esercizi sulle equazioni di secondo grado complete:
$$
3(x-1)(2x+2)=0.
$$
Quasi tutti faranno il $\Delta$ e tutta la ricetta. Poi torno indietro e gli chiedo se c'è una strategia migliore. Gli faccio vedere quanto è comodo.

Gliene faccio un'altra:
$$
(x+2)(x-\frac 12)=0.
$$

Gliene faccio una di quinto grado:
$$
x(2x+1)(3x-1)(4x^2-1)(x^2+2x+2)=0.
$$

Noto che, in generale, è molto bello scrivere le equazioni in questa forma. Questa forma si dice \textit{“fattorizzata”}.

Cosa vuol dire fattorizzazione? Scrivere un polinomio di grado $n$ come prodotto di polinomi di grado inferiore. Esempio: invece di $2x^2-4x-6$ scrivere $2(x-3)(x+1).$


Come si fa?

Se io scrivo
$$
(x-1)(x+2)=0
$$
so che le soluzioni sono $x=1$ o $x=-2$. 

e se io sapessi le soluzioni? Come faccio a fattorizzare? Se so che le soluzioni sono $x_1 = 2$ e $x_2=-1$? Allora potrebbe essere che l'eq di partenza sia 
$$
(x-2)(x+1)=0
$$

Ma potrebbe essere anche 
$$
2(x-2)(x+1)=0.
$$

Come faccio in generale?

\begin{mdframed}
    \textbf{Teorema} (della fattorizzazione).
    Presentarlo direttamente sotto forma di flowchart.
\end{mdframed}

\subsubsection{Esercizi sulla fattorizzazione}

Fattorizza, se possibile, i seguenti polinomi:
\begin{enumerate}
    \item $3x^2-12x+12$ $\rightarrow$ $\Delta=0$, una sola soluzione $x=2$: \quad $3(x-2)^2$
    \item $2x^2-4x-2$ $\rightarrow$ $\dfrac\Delta4=8$, soluzioni $x=1\pm\sqrt2$: \quad $2\left(x-1-\sqrt2\right)\left(x-1+\sqrt2\right)$
    \item $2x^2-3x+4$ $\rightarrow$ $\Delta=-23<0$, nessuna soluzione: non si può fattorizzare
    \item $2x^5-2x^4-12x^3$ $\rightarrow$ raccolgo $2x^3$: \quad $2x^3\left(x^2-x-6\right)=2x^3(x-3)(x+2)$
\end{enumerate}

Semplifica le seguenti frazioni, fattorizzando numeratore e denominatore:
\begin{enumerate}
    \setcounter{enumi}{4}
    \item $\dfrac{2x^2+x-1}{x^2+4x+3}$ $\rightarrow$ $\dfrac{(2x-1)(x+1)}{(x+1)(x+3)}=\dfrac{2x-1}{x+3}$, \quad con $x\neq-1$, $x\neq-3$
    \item $\dfrac{x^6-x^5-2x^4}{3x^3-6x^2}$ $\rightarrow$ raccolgo $x^4$ sopra e $3x^2$ sotto: \quad $\dfrac{x^4(x-2)(x+1)}{3x^2(x-2)}=\dfrac{x^2(x+1)}{3}$, \quad con $x\neq0$, $x\neq2$
\end{enumerate}

\section{Equazioni Frazionarie}

Consideriamo l'equazione $x^2=x+1$. Possiamo risolverla normalmente come un'equazione completa, seguendo la procedura del $\Delta$, e otterremmo $x_{1,2}=\frac{1\pm\sqrt 5}{2}.$ 

Però potremmo farlo anche in modo più creativo.
Se assumiamo che il valore di $x$ sia diverso da zero, possiamo dividere per $x$ e otteniamo

$$
x=1+\frac 1 x.
$$

Quindi abbiamo trovato che il valore di $x$ dipende da $x$ stesso. Siamo cioè in presenza di quello che in Matematica si chiama \textit{ricorsione}. Effettivamente possiamo sostituire la $x$ del membro di destra con l'espressione completa per $x$. Di questo modo, otteniamo

$$
x=1\frac{1}{1+\frac 1 x}.
$$
Possiamo ripetere questo gioco, anche infinite volte, e otteniamo (scrivere il castello di frazioni). Sappiamo che questo numero



\subsection{Scaletta: frazionarie con il denominatore da scomporre}

Due ore da $50$ minuti, con una pausa di $5$ minuti a metà. Le frazionarie con il denominatore di primo grado le hanno già fatte: qui l'unica cosa nuova è che \textbf{prima bisogna scomporre il denominatore}, e alla fine resta un'equazione di secondo grado. Un esercizio per ogni idea nuova, niente doppioni. Fra parentesi quadre le risposte.

\subsubsection*{Prima ora}

\paragraph{1. Riscaldamento: scomporre (0--8')}
Alla lavagna, uno per tipo di scomposizione che useremo:
\begin{enumerate}
    \item $x^2-9$ \quad [differenza di quadrati: $(x-3)(x+3)$]
    \item $x^2-4x$ \quad [raccoglimento: $x(x-4)$]
    \item $x^2-5x+6$ \quad [trinomio, con le soluzioni $2$ e $3$: $(x-2)(x-3)$]
\end{enumerate}

\paragraph{2. La ricetta (8--15')}
La scrivo in un angolo della lavagna e la lascio lì per tutte e due le ore.
\begin{mdframed}
    \begin{enumerate}
        \item \textbf{Scomponi} tutti i denominatori.
        \item Scrivi le \textbf{condizioni di esistenza}: ogni fattore del denominatore deve essere diverso da $0$.
        \item Trova il \textbf{denominatore comune} (il prodotto di tutti i fattori, presi una volta sola) e \textbf{moltiplica} tutti e due i membri per esso: i denominatori spariscono.
        \item \textbf{Risolvi} l'equazione di secondo grado che resta.
        \item \textbf{Controlla} le soluzioni con le condizioni di esistenza: quelle che non le rispettano si scartano.
    \end{enumerate}
\end{mdframed}

\paragraph{3. Insieme alla lavagna: differenza di quadrati (15--30')}
$$
\frac{x}{x-2}-\frac{1}{x+2}=\frac{8}{x^2-4}
$$
[$x^2-4=(x-2)(x+2)$; CE: $x\neq2$, $x\neq-2$. Moltiplico per $(x-2)(x+2)$: $x(x+2)-(x-2)=8$, cioè $x^2+x-6=0$, quindi $(x+3)(x-2)=0$. Le soluzioni sarebbero $x=-3$ e $x=2$, ma $x=2$ \textbf{non rispetta le CE}: resta solo $x=-3$.]

Qui mi fermo sul punto 5: senza le CE avrebbero tenuto anche $x=2$. Far sostituire $x=2$ nell'equazione di partenza: viene una divisione per zero.

\paragraph{4. Da soli: raccoglimento (30--50')}
$$
\frac{x}{x-1}+\frac{2}{x}=\frac{6}{x^2-x}
$$
[$x^2-x=x(x-1)$; CE: $x\neq0$, $x\neq1$. Moltiplico per $x(x-1)$: $x^2+2(x-1)=6$, cioè $x^2+2x-8=0$, quindi $(x+4)(x-2)=0$: $x=-4$ e $x=2$, \textbf{tutte e due accettabili}.]

Lascio un buon quarto d'ora per farlo davvero da soli, poi lo correggo alla lavagna con uno di loro. È quella in cui non si scarta niente: così vedono che non succede sempre.

\subsubsection*{Pausa (50--55')}

\subsubsection*{Seconda ora}

\paragraph{5. Da soli: trinomio (55--67')}
$$
\frac{x}{x-3}-\frac{1}{x-2}=\frac{1}{x^2-5x+6}
$$
[$x^2-5x+6=(x-2)(x-3)$; CE: $x\neq2$, $x\neq3$. Moltiplico: $x(x-2)-(x-3)=1$, cioè $x^2-3x+2=0$, quindi $(x-1)(x-2)=0$: $x=1$ accettabile, $x=2$ \textbf{da scartare}.]

Se si bloccano sulla scomposizione, ricordo che le soluzioni di $x^2-5x+6=0$ le abbiamo trovate nel riscaldamento.

\paragraph{6. Attenzione al meno (67--76')}
$$
\frac{2x}{x-1}-\frac{x+3}{x+1}=\frac{4}{x^2-1}
$$
[CE: $x\neq\pm1$. Moltiplico per $(x-1)(x+1)$: $2x(x+1)-(x+3)(x-1)=4$. Il meno davanti alla seconda frazione cambia il segno di \textbf{tutto} il prodotto: $2x^2+2x-(x^2+2x-3)=4$, cioè $x^2-1=0$, quindi $x=\pm1$: \textbf{tutte e due da scartare}, l'equazione è impossibile.]

Prima di correggere, guardo chi ha scritto $-x^2+2x-3$: è l'errore più frequente, e conviene farlo vedere alla lavagna. Il risultato \lette{impossibile} li spiazzerà: è voluto.

\paragraph{7. Due scomposizioni di seguito (76--85')}
$$
\frac{x}{x-2}-\frac{4}{x+2}=\frac{32}{2x^2-8}
$$
[$2x^2-8=2(x^2-4)=2(x-2)(x+2)$: prima raccoglimento, poi differenza di quadrati. CE: $x\neq\pm2$. Moltiplico per $2(x-2)(x+2)$: $2x(x+2)-8(x-2)=32$, cioè $2x^2-4x-16=0$, e dividendo per $2$: $x^2-2x-8=0$, quindi $(x-4)(x+2)=0$: $x=4$ accettabile, $x=-2$ da scartare.]

Novità: il $2$ del denominatore. Se qualcuno lo dimentica nel denominatore comune, i conti non tornano: farlo notare.

\paragraph{8. Una storiella per chiudere (85--97')}
Un gruppo di amici affitta una casa al mare per $600$ euro, divisi in parti uguali. Se fossero in $2$ di più, ognuno pagherebbe $25$ euro in meno. Quanti sono?

[Se sono $x$, ognuno paga $\frac{600}{x}$; in $x+2$ pagherebbero $\frac{600}{x+2}$ a testa. Equazione: $\dfrac{600}{x}-\dfrac{600}{x+2}=25$. CE: $x\neq0$, $x\neq-2$. Moltiplico per $x(x+2)$: $600(x+2)-600x=25x(x+2)$, cioè $1200=25x^2+50x$; divido per $25$: $x^2+2x-48=0$, quindi $(x+8)(x-6)=0$. Le CE vanno bene per tutte e due, ma $x=-8$ non ha senso: \textbf{sono $6$ amici}.]

Qui il denominatore non va scomposto: la novità è che, oltre alle CE, alla fine si scarta anche la soluzione che non ha senso nel problema. Imposto io l'equazione insieme a loro, i conti li fanno loro.

\paragraph{Compiti (97--100')}
\begin{enumerate}
    \item $\dfrac{x+2}{x}+\dfrac{1}{x+1}=\dfrac{7}{x^2+x}$ \quad [CE: $x\neq0$, $x\neq-1$; $x^2+4x-5=0$: $x=-5$, $x=1$, tutte e due accettabili]
    \item $\dfrac{x}{x+3}+\dfrac{3}{x-3}=\dfrac{18}{x^2-9}$ \quad [CE: $x\neq\pm3$; resta $x^2-9=0$, cioè $x=\pm3$: tutte e due da scartare, \textbf{impossibile}]
    \item $\dfrac{x}{x-2}+\dfrac{1}{x+1}=\dfrac{9}{3x^2-3x-6}$ \quad [$3x^2-3x-6=3(x-2)(x+1)$: raccoglimento, poi trinomio. CE: $x\neq2$, $x\neq-1$. Moltiplico per $3(x-2)(x+1)$: $3x(x+1)+3(x-2)=9$, cioè $x^2+2x-5=0$: $x=-1\pm\sqrt6$, tutte e due accettabili]
    \item $\dfrac{x}{x-3}=\dfrac{2}{x-2}+\dfrac{2}{x^2-5x+6}$ \quad [$x^2-4x+4=0$: $x=2$, da scartare: \textbf{impossibile}]
    \item Una gita in pullman costa in tutto $900$ euro, divisi fra i partecipanti. Se ci fossero $5$ partecipanti in più, ognuno pagherebbe $6$ euro in meno. Quanti sono i partecipanti? \quad [$\frac{900}{x}-\frac{900}{x+5}=6$, cioè $x^2+5x-750=0$, $\Delta=3025=55^2$: $x=25$]
\end{enumerate}

\paragraph{Se resto indietro}
Nell'ordine: il punto 7 (lo do per casa), poi la storiella (la imposto io e la finiscono a casa). \textbf{Non} tagliare il punto 3, che fa vedere perché servono le CE, né il punto 6, che contiene l'errore di segno più frequente.

\section{Note sul capitolo}
Alla sezione delle equazioni monomie, hanno avuto abbastanza difficoltà a comprendere che $-2(x-3)^2=0$ aveva come soluzione $x=3$. Ho provato a farlo fare a loro e ci è riuscito Guglielmo ma solo dopo un po' di tempo. Ha proposto di dividere per $-2$ e da lì hanno compreso.








\end{document}